\documentclass[11pt]{article}
\usepackage[margin=1in]{geometry}
\usepackage{amsmath,amssymb}
\usepackage{enumitem}
\begin{document}

\section*{Physics-Informed Forward Soil EC Model}

The current forward model predicts electrical conductivity from soil-health
state variables and soil context variables.

\subsection*{Response}
\[
y = \log_{10}(\mathrm{EC})
\]

\subsection*{Model Form}
\[
\hat{y}
= \beta_0
+ \sum_i \beta_i^{(s)} \, g_i(s_i)
+ \sum_j \beta_j^{(c)} \, h_j(c_j)
+ \sum_{(i,j)\in\mathcal{I}} \beta_{ij}^{(sc)} \, g_i(s_i)\,h_j(c_j)
\]

\[
\widehat{\mathrm{EC}} = 10^{\hat{y}}
\]

\subsection*{State Variables}
\[
\mathbf{s} =
\left\{
N,\;P,\;K,\;\mathrm{CEC},\;\mathrm{OC},\;\mathrm{Moisture}
\right\}
\]

At present, moisture is included only if available in the data.

\subsection*{Context Variables}
\[
\mathbf{c} =
\left\{
\mathrm{Clay},\;\mathrm{Silt},\;\mathrm{Sand},\;\mathrm{Coarse},\;
\mathrm{pH}_{\mathrm{CaCl_2}},\;\mathrm{pH}_{\mathrm{H_2O}},\;\mathrm{CaCO_3}
\right\}
\]

\subsection*{Transforms}
For skewed positive variables, the model uses:
\[
g(x)=\log(1+x),
\qquad
h(x)=\log(1+x)
\]

For variables that are not strongly skewed, the model uses the raw value:
\[
g(x)=x
\qquad \text{or} \qquad h(x)=x
\]

\subsection*{Selected Interaction Set}
The current interaction set couples each state variable with the key context variables:
\[
\mathcal{I} =
\left\{
\mathrm{Clay},\;\mathrm{Silt},\;\mathrm{Sand},\;
\mathrm{pH}_{\mathrm{H_2O}},\;\mathrm{pH}_{\mathrm{CaCl_2}},\;\mathrm{CaCO_3}
\right\}
\]

So, for example, the model includes terms such as:
\[
g(N)\,h(\mathrm{Clay}),\quad
g(K)\,h(\mathrm{pH}_{\mathrm{H_2O}}),\quad
g(\mathrm{OC})\,h(\mathrm{CaCO_3})
\]

\subsection*{Interpretation}
\begin{itemize}[leftmargin=1.5em]
\item State variables represent soil-health quantities of interest.
\item Context variables condition how those state variables translate into conductivity.
\item Clay is treated as a context variable because it modifies the EC response rather than acting only as a health target.
\end{itemize}

\end{document}
